% 原书物理页=印刷页21；3.1 全节.
\chapter{有限群的结构}\label{chap:fg-3}
\section{Sylow 定理}\label{sec:fg-3-1}

在本节中，$p$ 恒表示一个素数，$G$ 表示一个有限群. 称 $G$ 是一个\newterm{fg-p-group}{$p$-群}{$p$-group}，如果它的阶是 $p$ 的幂 $p^n$（其中 $n\geqslant0$）.

\begin{definition}\label{def:fg-3-1-1}
称 $G$ 的子群 $S$ 是 $G$ 的一个 \newterm{fg-sylow-subgroup}{Sylow $p$-子群}{Sylow $p$-subgroup}（或更简单地：$G$ 的一个 $p$-Sylow），如果 $S$ 是一个 $p$-群，且其指数 $[G:S]$ 与 $p$ 互素.
\end{definition}

换言之，若把 $|G|$ 写成 $|G|=p^nm$，其中 $m$ 与 $p$ 互素，则 $G$ 的一个 $p$-Sylow 就是一个阶为 $p^n$ 的子群. 此时显然，与 $S$ 共轭的每个子群仍是 $G$ 的 $p$-Sylow. 而不那么显然的则是：$p$-Sylow 是否存在！

\begin{example}\label{ex:fg-3-1-2}
设 $k$ 是一个有限域，其基数为 $q=p^f$. 群 $G=\mathrm{GL}_n(k)$ 的阶为 $p^{fn(n-1)/2}m$，其中 $m=\prod_{i=1,\dots,n}(q^i-1)$ 与 $p$ 互素. $G$ 中由严格上三角矩阵（即对角线上只有 $1$ 的矩阵）组成的子群 $S$ 是 $G$ 的一个 $p$-Sylow.
\end{example}

\begin{theorem}\label{thm:fg-3-1-3}
设 $G$ 是有限群，$p$ 是一个素数. 那么：
\begin{enumerate}[label=\roman*)]
\item $G$ 拥有 $p$-Sylow；
\item $G$ 的任意两个 $p$-Sylow 相互共轭；
\item $G$ 的任一 $p$-子群都包含在 $G$ 的某个 $p$-Sylow 中.
\end{enumerate}
\end{theorem}

存在性（第~i)~点）的证明在于借助下面的引理化归为一个显式的情形：

\begin{lemma}\label{lemma:fg-3-1-4}
设 $G$ 是有限群，$H$ 是 $G$ 的一个子群，$S$ 是 $G$ 的一个 $p$-Sylow. 那么存在 $g\in G$，使得 $H\cap gSg^{-1}$ 是 $H$ 的一个 $p$-Sylow.
\end{lemma}

\begin{proof}
令 $X=G/S$ 为 $G$ 模 $S$ 的左陪集全体构成的集合，群 $H$ 通过左平移作用在其上. 在 $H$ 的作用下，$X$ 中任一点 $gS$ 的固定子是 $H$ 与 $S$ 的共轭子群 $gSg^{-1}$ 的交，因而是一个 $p$-群. 由于 $p$ 不整除 $\operatorname{card}(X)=[G:S]$，$X$ 在 $H$ 作用下的各条轨道中至少有一条，设为 $H.gS$，其基数 $[H:(H\cap gSg^{-1})]$ 与 $p$ 互素. 故 $H\cap gSg^{-1}$ 是 $H$ 的一个 $p$-子群，其阶与其在 $H$ 中的指数互素，也就是说它是 $H$ 的一个 $p$-Sylow.
\end{proof}

\begin{proof}[定理的证明]
i) 根据引理~\ref{lemma:fg-3-1-4}\footnote{译注：原文此处作``Lemme 1''，系原书交叉引用失效，即引理~\ref{lemma:fg-3-1-4}.}，只须把 $G$（也就是构造 $G$ 的一个单同态）嵌入到一个拥有 $p$-Sylow 的群 $\mathbf{G}$ 中. 而若 $n=|G|$，则 $G$ 可嵌入对称群 $S_n$，后者又通过同态 $\rho$ 嵌入 $\mathbf{G}=\mathrm{GL}_n(\mathbb{F}_p)$：$\rho$ 给每个 $\sigma\in S_n$ 附上 $\mathbb{F}_p$-向量空间 $\mathbb{F}_p^n$ 的自同构 $\rho(\sigma)$，其定义为 $\rho(\sigma)(x_1,\dots,x_n)=(x_{\sigma^{-1}(1)},\dots,x_{\sigma^{-1}(n)})$. 根据定理上方的例子，$\mathbf{G}$ 拥有 $p$-Sylow，故 $G$ 亦然.
ii)（相应地，iii)）把引理应用于 $H$ 是 $G$ 的一个 $p$-Sylow（相应地，一个 $p$-子群）的情形.
\end{proof}

\begin{corollary}\label{cor:fg-3-1-5}
$G$ 总是一个有限群.\footnote{译注：原文如此，此句疑为原书笔误或排版残留；按 i)、ii) 的内容，此处本应是对 $G$ 有限的假定（``设 $G$ 为有限群''）.}
\begin{enumerate}[label=\roman*)]
\item 对任意整除 $|G|$ 的素数 $p$，$G$ 拥有阶为 $p$ 的元素；
\item $G$ 是 $p$-群当且仅当它的所有元素的阶都是 $p$ 的幂.
\end{enumerate}
\end{corollary}

\begin{proposition}\label{prop:fg-3-1-6}
若 $|G|=p^nm$，其中 $(p,m)=1$，则 $G$ 的 Sylow $p$-子群的个数 $s=\operatorname{card}(\mathcal{S})$ 整除 $m$，且满足 $s\equiv1\bmod p$.
\end{proposition}

\begin{proof}
根据定理~\ref{thm:fg-3-1-3}\footnote{译注：原文此处仅作``D'après le Théorème''而未标编号，依上下文即定理~\ref{thm:fg-3-1-3}~ii).}，$G$ 通过共轭作用在 $\mathcal{S}$ 上，且该作用是传递的. 而 $\mathcal{S}$ 中元素 $S$ 在 $G$ 中的固定子是 $N_G(S)$，它包含 $S$. 因此 $s=[G:N_G(S)]$，故它整除 $[G:S]=m$. 现在让 $S$ 通过共轭作用在 $\mathcal{S}$ 上，并证明唯一的一元轨道就是 $S$ 自身所在的轨道：若对一切 $s\in S$ 都有 $sS's^{-1}=S'$，则 $SS'$ 是 $G$ 的子群，其所有元素的阶都是 $p$ 的幂；根据推论~\ref{cor:fg-3-1-5}~ii)\footnote{译注：原文此处作``corollaire 1.ii''，系原书交叉引用失效，即推论~\ref{cor:fg-3-1-5}~ii).}，它是一个 $p$-群，于是由 $p$-Sylow 的极大性得 $SS'=S=S'$. 其余轨道的基数不等于 $1$ 且整除 $|S|$，因而被 $p$ 整除. 于是 $s\equiv1\bmod p$.
\end{proof}
